% Notes or not
%%\documentclass[10pt,t,handout]{beamer}
%%%%
\documentclass[10pt]{beamer}
%%
\usepackage{pgfpages}
\usepackage{graphicx}
\usepackage{ulem}
\usepackage{color}
\usepackage{fancyvrb}
% get rid of junk
\usetheme{default}
\beamertemplatenavigationsymbolsempty
\hypersetup{pdfpagemode=UseNone} % don't show bookmarks on initial view
% font
\usefonttheme{professionalfonts}
\usefonttheme{serif}
% page number
\setbeamertemplate{footline}{%
    \raisebox{5pt}{\makebox[\paperwidth]{\hfill\makebox[20pt]{\color{gray}
          \scriptsize\insertframenumber}}}\hspace*{5pt}}
% add a bit of space at the top of the notes page
  \addtobeamertemplate{note page}{\setlength{\parskip}{12pt}}

\setbeamertemplate{blocks}[rounded][shadow=true]

% Notes or not
\setbeameroption{show notes}
%%\setbeamertemplate{note page}[plain]
%%\pgfpagesuselayout{2 on 1}
%%%%
%%\setbeameroption{hide notes}

\newenvironment{xframe}[2][]
  {\begin{frame}[fragile,environment=xframe,#1]
  \frametitle{#2}}
  {\end{frame}}

\title{Workshop: R and Bioinformatics}
\author{Jean Monlong \& Simon Papillon}
\institute{Human Genetics department}
\date{October 28, 2013}

\begin{document}

%%%%%%%%%%%%%%%%%%%%
%% Title Slide
\begin{frame}
  \titlepage
  \centering
  \includegraphics[page=1,height=.1\textheight]{imgs/McGill-Logo1.png}

\end{frame}

%%%%
%%%%%%%%%%%%%%%%%%%%%%%
\section{Why using R for bioinformatics ?}
%%%%%%%%%%%%%%%%%%%%%%%
%%%%

\begin{frame}[label=handout]{Why using R for bioinformatics ?}
  \begin{itemize}
  \item Flexible statistics and data visualization software.
  \item Many packages and a vast community: Bioconductor.
  \item Simple and easy, compared at other computing languages.
  \end{itemize}
  \note{671 packages in Bioconductor. Bioconductor provides tools for the analysis and comprehension of high-throughput genomic data.}
\end{frame}

%%%%%%%

\begin{frame}[label=handout]{Today's workshop}
  \begin{block}{Goal}
    \begin{itemize}
    \item Explain and demonstrate some key principles to do good bioinformatics.
    \item How to structure the analysis, how to write efficient script, explore your data. 
    \item What to do and not to do...
      \bigskip
    \item This is {\bf NOT} a package tutorial but will point at useful resources.
    \end{itemize}    
  \end{block}
  \medskip

  \begin{block}{Learning from your mistakes}
    You will get errors today, before raising your hand:
    \begin{itemize}
    \item Check your command for typos.
    \item Try to understand the error message.
    \item Check your input/objects. 
    \end{itemize}    
  \end{block}
\end{frame}

%%%%%%%


%%%%%%%

%%%%%%%
\section{Reminder and updates}

\begin{xframe}{Functions}
  {\it\small Previously on HGSS workshops: }
  \begin{block}{}
    \begin{description}
      \item[function] To define functions.
      \item[return] Define what will be returned by the function. 
    \end{description}
    All the object created within the function are temporary.
    \bigskip
    \end{block}
    \begin{block}{Structure}    
\begin{Verbatim}[commandchars=\\\{\}]
\color{green!60!black}myFunctionName \color{blue}<- function(\color{green!60!black}input.obj1\color{blue},\color{green!60!black}second.input.obj \color{blue}) \{
\color{black}...
... Intructions on 'input.obj1' and 'second.input.obj'
...
\color{blue}return(\color{green!60!black}my.output.obj\color{blue})
\color{blue}\}

\color{black}myFunctionName(1,c(2,4,5))
\end{Verbatim}
  \end{block}
\end{xframe}

\begin{xframe}{Conditions}
  \begin{block}{Logical tests}
	
    \begin{description}
    \item[{\sf ==}] Are both values equal.
    \item[{\sf $>$ or $>=$}] Is the left value greater (greater or equal)
    than the right value.
	\item[{\sf $<$ or $<=$}] Is the left value smaller (smaller or equal) than
	the left value.
    \item[!] Is a NOT operator that negates the value of a test.
    \item[$|$] Is an OR operator used to combine logical tests. Returns TRUE if
    either are TRUE.
    \item[\&] Is an AND operator used to combine logical tests. Returns TRUE
    if both are TRUE
    \end{description}
  \end{block}
  \begin{exampleblock}{Example}
\begin{verbatim}
test <- 2 + 2 == 4    ## (TRUE)
!test                 ## (FALSE)
test | !test          ## (TRUE)
test & !test          ## (FALSE)
\end{verbatim}  
  \end{exampleblock}
  \note{Here I changed the $||$ and \&\& by the simple $|$ and \&. I know they should be used for vectorized test but it works and it's simpler for them to remember.}
\end{xframe}

%%%%%%%%%%%

\begin{xframe}{Conditions}
  \begin{block}{Boolean}
  Any logical tests can be vectorized (compare 2 {\sf vector}s).
    \begin{description}
    %\item[$\mid$] Is a OR operator for vectorized application.
    %\item[\&] Is an AND operator for vectorized application.
    \item[which] Returns the index of the {\sf vector}s with {\it TRUE} values.
    %%\item[any] Take a {\sf vector} of {\it logical} and return {\it TRUE} if at least one value is {\it TRUE}.
    %%\item[\%in\%] Vectorized any. See example/supp material.
    \end{description}
  \end{block}
  \begin{exampleblock}{Example}
\begin{verbatim}
5:8 == 6                         ## FALSE,TRUE,FALSE,FALSE
5:8 >= 6 & 5:8<=7                ## FALSE,TRUE,TRUE,FALSE

c(TRUE, TRUE) & c(TRUE, FALSE)   ## TRUE,FALSE
c(TRUE, FALSE) | c(FALSE, FALSE) ## TRUE,FALSE

which(5:10 == 6)                 ## 2
which(5:10 > 6)                  ## 3,4,5,6
\end{verbatim}  
  \end{exampleblock}
\end{xframe}
%%%%%%%%%

\begin{frame}{Conditions - Exercise}
  \begin{block}{Exercise}
  Create a function that: 
  \begin{enumerate}
  \item remove values below $3$ from a {\sf vector}.
  \item  remove values below a specified threshold from a {\sf vector}.
  \end{enumerate}
  \end{block}

  \bigskip

  \begin{block}{For more advanced users}
    Have a look at these functions on {\it logical} vectors:
    \begin{itemize}
    \item any, $\%$ in$\%$.
    \item sum, mean, table.
    \end{itemize}
  \end{block}

  \bigskip

  \begin{block}{Extra tips}
    \begin{itemize}
    \item Don't filter out, keep in.
    \item Use the boolean vector directly between $[~]$.
    \end{itemize}
  \end{block}
\end{frame}

%%%%%%%

%%%%%%%

\begin{xframe}{Testing conditions}
  \begin{block}{{\sf if else}}
    Test a condition, if {\it TRUE} run some instruction, if {\it FALSE} something else (or nothing).
\begin{verbatim}
if( Condition ){
...   Instructions
} 
\end{verbatim}  
  \end{block}
  \begin{exampleblock}{Example}
\begin{verbatim}
if(length(luckyNumbers)>3){
  cat("Too many lucky numbers.\n")
  luckyNumbers = luckyNumbers[1:3]
} else if(length(luckyNumbers)==3){
  cat("Just enough lucky numbers.\n")
} else {
  cat("You need more lucky numbers.\n")
}
\end{verbatim}  
  \end{exampleblock}
\end{xframe}

\begin{xframe}{Loops}
  \begin{block}{{\sf for} loops}
    Iterate over the element of a container and run instructions.
\begin{verbatim}
for(v in vec){
...  Instruction
}
\end{verbatim}  
  \end{block}
  \begin{block}{{\sf while} loops}
    Run instructions as long as a condition is {\it TRUE}.
\begin{verbatim}
while( CONDITION ){
...  Instruction
}
\end{verbatim}  
  \end{block}
  \begin{exampleblock}{Example}
\begin{verbatim}
facto = 1
for(n in 1:10){
   facto = facto * n
}
\end{verbatim}  
  \end{exampleblock}
\end{xframe}

%%%%%%%%%

\begin{frame}{Exercises}
  \begin{block}{{\sf if else}}
    Create a function that classify the average value of a {\sf vector}. It returns:
    \begin{itemize}
    \item {\it low} if the average if below $3$.
    \item {\it medium} if the average if between $3$ and $7$.
    \item {\it high} if the average if above $7$.
    \end{itemize}
  \end{block}
  
  \bigskip
  
  \begin{block}{Loops}
    Write a function that computes the mean values of a matrix columns:
    \begin{enumerate}
    \item using the {\sf apply}  function.
    \item using a {\sf for} loop.
    \item (using a {\sf while} loop.)
    \end{enumerate}
  \end{block}
\end{frame}


%%%%%%%%%

\begin{frame}{Important principles}
  \begin{block}{Scripting}
    Write scripts of your analysis: 
    \begin{itemize}
    \item Keeping track, easy rerun, easy parameter tweaking.
    \item {\sf Rstudio} or other interfaces ({\sf Emacs+ESS},...).
    \end{itemize}
  \end{block}
  \begin{block}{Clear and modular code}
    \begin{itemize}
    \item Define clear analysis steps.
    \item Write function(s) for each step.
      \begin{itemize}
      \item Keeps the data and parameters used clear (for you and for R).
      \item No confusing temporary objects.
      \item No repeating code.
      \item More suitable for {\sf apply}-like usage.
      \item Easy parameter tweaking.
      \end{itemize}
    \end{itemize}
  \end{block}
  \begin{block}{Efficiency matters}
    \begin{itemize}
    \item Data structure and manipulation.
    \item Especially relevant with our large data.
    \end{itemize}

  \end{block}
  \note{Simon what interface do you use ?}
\end{frame}

%%%%%%%%%

\begin{frame}{Data exploration - All you can plot}
  \centering
  \uline{{\bf Answer your questions with plots !}}
  \medskip

  \begin{block}{Utility}
    \begin{itemize}
    \item Get an idea of the quality of the data and potential issues.
    \item Get a full answer.
    \item Detect potential biases.
    \item (Find unexpected results.)
    \end{itemize}
  \end{block}
  
  \begin{block}{Through all your analysis}
    \begin{itemize}
    \item Quality Control plots at the beginning.
    \item Control plot after each steps.
    \item Awesome plot with your results.
    \end{itemize}
  \end{block}
\end{frame}

%%%%%%%%%

\begin{xframe}[shrink=5]{Today's dataset - Methylation analysis}
  \begin{block}{Goal}
    Analyze methylation data following those principles. 
  \end{block}
  \begin{block}{DNA methylation data}
    \begin{itemize}
    \item Identify sample outliers (and remove them).
    \item Identify co-variates (sex, age) using PCA.
    \item Use heatmap to plot sample groupings.
    \item Point out differentially methylated sites.
    \item Plot methylation levels of interesting sites.
    \end{itemize}
  \end{block}
  \begin{exampleblock}{The data set}
\begin{verbatim}
beta_value # The methylation data
probe_data # The annotation of each probe
pheno_data # The annotation of each sample (metadata)
\end{verbatim}  
  \end{exampleblock}
  \begin{exampleblock}{Get to know your data}
    Check the first few rows of each object (remember the {\sf head} function). Check the type and size of your data({\sf str}).
  \end{exampleblock}
\end{xframe}

%%%%%%%

\begin{frame}{Importing and manipulating large files}
  Extra parameters for read.table: nb of lines, column classes.
  \note{Not sure where to put this. If your data is not too big, maybe we can talk about it in the end, but it would be nice to tell them what not to do in the beginning.}
\end{frame}

%%%%%%%

\begin{xframe}{Application - Identify outliers}
  \begin{block}{Checking the distribution}
    \begin{itemize}
      \item Identify the {\sf case} samples. 
      \item Take a look at the density of those samples.
      \item Are there any samples that stand out ?
    \end{itemize}
  \end{block}
  \begin{exampleblock}{Useful functions}
\begin{verbatim}
density(x) # Compute the density of x
plot(x) # Create a plot of x
plot(density(x)) # Plot the density of x
\end{verbatim}
  \end{exampleblock}
  \note{Distribution of what? methylation ?}
\end{xframe}

%%%%%%%

\begin{frame}{Quality Control plots}
  \begin{block}{Aim}
    Assess the quality of your data and potential artefacts that could bias your analysis.
  \end{block}
  \begin{block}{Basic approaches}
    \begin{itemize}
    \item Principal Component Analysis: representing the largest variation in the data.
    \item Clustering: summarizing similarity relations between samples/genes.
    \item Heatmaps: 
    \item Testing metadata: gender, age, ...
    \end{itemize}
  \end{block}
  \note{wikipedia cartoon for PCA, cluster tree\\How to define simply heatmaps and its utility...}
\end{frame}

%%%%%%%%%

\begin{xframe}[shrink=10]{Quality Control plots - Functions}
  \begin{block}{PCA using {\sf prcomp}}
    PCA of the matrix columns; plot of the variance explained by the first PCs; representation of the {\bf rows} using the first two PCs.
\begin{verbatim}
peeSeaAye = prcomp(input.matrix)
plot(peeSeaAye)
plot(peeSeaAye$x,type="n")
text(peeSeaAye$x,labels=rownames(input.matrix))
\end{verbatim}  
  \end{block}
  \begin{block}{Clustering using {\sf hclust}}
    Clustering using a \uline{distance} matrix, e.g. from correlation between {\bf columns}.
\begin{verbatim}
cor.dist = as.dist(1-cor(input.matrix))

kleusteur = hclust(cor.dist,method="ward")
plot(kleusteur)

library(MASS)
mDeeS = isoMDS(cor.dist)
plot(mDeeS$points)
\end{verbatim}  
  \end{block}
  \note{highlight the functions\\Importance of the distance demonstration maybe...}
\end{xframe}

%%%%%%%%%%

\begin{frame}{Heatmaps}
  
\end{frame}

%%%%%%%%%%

\begin{xframe}{Application - Identify outliers}
  \begin{block}{Using the PCA}
    \begin{itemize}
      \item Using all samples, plot the first 2 Principle Components.
      \item Color the samples according to their Sample\_Group status. 
      \item How do the outliers identified previously behave ?
      \item Remove those samples from your data set (dont forget to propagate
      the changes in all objects!).
    \end{itemize}
  \end{block}
   \begin{exampleblock}{PCA example}
\begin{verbatim}
pca <- prcomp(t(beta_value))
plot(pca$x)
\end{verbatim}  
  \end{exampleblock}
\end{xframe}

%%%%%%%

%%%%%%%
\begin{xframe}{Application - Identify phenotypic co-variates}
  \begin{block}{Using the PCA}
    \begin{itemize}
      \item Re-compute the PCA without the outlier samples.
      \item Plot the 1st PC with the age of the samples (do you notice
      something ?)
      \item Color the sample according to their Sample\_Group status
      \item Re-compute the PCA using {\sf control} samples only.
	  \item Does sex have an effect on data ?
	  \item Knowing that females have 2 X chromosomes, 1 silenced by methylation
	  are you able to predict the sex of the samples?
	  \item Check your prediction in the PCA plot
    \end{itemize}
      \end{block}

\end{xframe}

%%%%%%%
\begin{xframe}[shrink=5]{Application - Predicting the sex}
  \begin{exampleblock}{Some tips}
\begin{verbatim}
# Here we assume that the X chromosome is more methylated
# in females
x_probes <- probe_data$chrom == "chrX"
summary(x_probes)

controls <- pheno_data$Sample_Group == "control"
summary(controls)
pca <- prcomp(t(beta_value[x_probes, controls]))
mean_x_meth <- apply(beta_value[x_probes, controls], 2, mean)
color_vec <- ifelse(mean_x_meth >= 0.45, "red", "blue")
plot(pca$x, col=color_vec)
\end{verbatim}  
  \end{exampleblock}
\end{xframe}

%%%%%%%

\begin{xframe}{Automated approach to identify metadata co-variates}
  \begin{block}{Linear regression}
    For example, between Principal Component and metadata. 
    \begin{itemize}
    \item Can be automated to test other Principal Components and numerous metadata. 
    \item To make a decision on borderline/unclear cases.
    \end{itemize}    
  \end{block}
  \begin{exampleblock}{Application}
\begin{verbatim}
summary(lm(pca$x[,1]~age))

covar.test = summary(lm(pca$x~age))
lapply(covar.test,function(l)l$coefficients[2,4])
\end{verbatim}  
  \end{exampleblock}
  \note{I can finish this once I see the data}
\end{xframe}


%%%%%%%
\begin{xframe}{Cluster the samples}
  \begin{block}{Using the heatmap}
    \begin{itemize}
      \item We dont need to use the full data set (too big)
      \item Let's find the 1,000 most variant sites and use them for clustering
      \item Make the heatmap (see previous examples)
      \item Color the sample according to their Sample\_Group status (use
      ColSideColors=color\_vec in heatmap)
    \end{itemize}
      \end{block}
    \begin{exampleblock}{Predicting the sex}
\begin{verbatim}
# Getting the 1,000 most variant sites
#
# Compute the variance
probe_var <- apply(beta_value, 1, var)
# Order them decreasingly
probe_var <- order(probe_var, decreasing = T)
# Get the top 1,000
probe_var <- probe_var[1:1000]
\end{verbatim}  
  \end{exampleblock}   

\end{xframe}

%%%%%%%

\begin{xframe}{Find differentially methylated sites}
  \begin{block}{A t-test compares 2 distributions}
    \begin{itemize}
      \item We want to compare {\sf cases} vs {\sf controls} at each site
      (loop)
      \item First, make a boolean vector for cases and one for controls (we saw
      how in previous slides)
    \end{itemize}
      \end{block}
    \begin{exampleblock}{t-test, an example}
\begin{verbatim}
t.test(rnorm(100), rnorm(100))

# Getting only the p-value
t.test(rnorm(100), rnorm(100))$p.value

\end{verbatim}  
  \end{exampleblock}   
    \begin{alertblock}{Non-parametric alternative}
      Wilcoxon test using {\sf wiclox.test} is a rank based test.
    \end{alertblock}
\end{xframe}

%%%%%%%%%

\begin{xframe}{Plot differentially methylated sites}
  \begin{block}{Using boxplots}
    \begin{itemize}
      \item Using a previously identified differentially methylated site make a
      boxplot of this site for {\sf cases} and {\sf controls}.
    \end{itemize}
  \end{block}
  
  AN EXAMPLE OF WHAT PLOT THEY SHOULD AIM AT
  
\end{xframe}

%%%%%%%


\begin{frame}{GenomicRanges and specific annotation}
  
\end{frame}

%%%%%

\begin{frame}{Extra: {\sf ggplot2}}
  \begin{block}{Introduction}
    A package to constuct beautiful and/or complex graphs. Many aspects of the graph are arranged automatically but everything can be specified. Easy layers addition.
  \end{block}
  
  EXAMPLE

\end{frame}

\begin{xframe}[shrink=10]{Extra: {\sf ggplot2} - {\sf data.frame} only}
    \begin{block}{{\sf data.frame}}
    The input object is always a {\sf data.frame}, with each rows being one {\it "point"} to represent and each column the different information on it.
    \medskip
    
    {\small {\sf data.frame}: in its simple version, a {\sf matrix} with different data type possible in each column.}
  \end{block}

  \begin{block}{Useful functions}
    \begin{description}
    \item[data.frame] To create a {\sf data.frame}.
    \item[subset] To subset a {\sf data.frame} using condition on the columns.
    \item[melt/reshape] To deconstruct a matrix into data.frame, or the opposite. {\it reshape} package.
    \item[aggregate] To compute summary statistics on subset of the data.frame.
    \item[ddply] {\sf apply}-like function on subset of the data.frame. {\it plyr} package.
    \end{description}
  \end{block}
  \begin{exampleblock}{Example}
\begin{verbatim}
gene.expression.df = data.frame(gene=c("A","B","C"),
                                   gene.expression=1:3)
dim(gene.expression.df)
gene.expression.df = subset(gene.expression.df, gene.expression>1)
\end{verbatim}  
  \end{exampleblock}
\end{xframe}

\begin{xframe}{Extra: {\sf ggplot2} - Simple graph}
  \begin{block}{}
    
  \end{block}
\end{xframe}

\section{More resource online}

\begin{frame}[shrink=10]{Online tutorials}
  \begin{block}{R}
    \begin{itemize}
    \item \url{http://www.twotorials.com/} : small video-tutorials.
    \item \url{http://www.r-tutor.com/} : R and statistics small web-tutorials.
    \item \url{www.youtube.com/user/rdpeng/} : Coursera {\it Computing for Data Analysis} videos. Other interesting videos, e.g. {\it ggplot2}.
    \item \url{http://cran.r-project.org/manuals.html} : R manual.
    \end{itemize}
  \end{block}
  \begin{block}{R and Bioinformatics}
    \begin{itemize}
    \item \url{http://stephenturner.us/p/edu} List of online resources for Bioinformatics.
    \item \url{http://bioinformatics.ca/workshops/2013/} : Bioinformatics workshop material.
    \item \url{http://manuals.bioinformatics.ucr.edu/home/R_BioCondManual} : Pieces of code for bioinformatics analysis, plots. Including Bioconductor.
    \item \url{http://bioconductor.org/help/course-materials/2013/} : Bioinformatics tutorials material: pdf and R scripts.
    \end{itemize}
  \end{block}
\end{frame}

\begin{frame}{Useful packages}
  \begin{description}
  \item[GenomicRanges] 
  \item[AnnotationHub]
  \item[TxBD]
  \end{description}
\end{frame}

\section{Extra slides}

%%%%%%%

\begin{xframe}{Lists}
  \begin{block}{Flexible container}
    A {\sf list} can contain any element type. It does not require elements to be of
    the same type.
    \begin{description}
      \item[list] Create a {\sf list}.
      \item[{l[[i]]} ] Get or set the $i^{th}$ object of the {\sf list}.
      \item[l\$toto] Get or set the element labeled as {\it toto}.
      \item[names] Get or set the names of the {\sf list} elements.
      \item[length] Get the number of element in the {\sf list}.
      \item[str] Output the structure of a R object.
    \end{description}
  \end{block}
  \begin{exampleblock}{Example}
\begin{verbatim}
l = list(vec=1:10,mat=matrix(runif(25),5))
str(l)
l
l$vec = 1
l
\end{verbatim}
  \end{exampleblock}
  \note{Questions:\\Make a phonebook: A list of 3 elements (vectors): names,
  phone number and address}
\end{xframe}

%%%%%%%

\begin{xframe}{Functions - {\sf lapply}}
  \begin{block}{apply for {\sf list}s}
    \begin{itemize}
    \item Useful way to iterate through {\sf list}s.
    \end{itemize}
  \end{block}
  \begin{exampleblock}{Example}
\begin{verbatim}
file_list <- list.files('.')
files_content <- lapply(file_list, function(file) \{
	data <- read.csv(file)
	#Do something with the data
	return(data)
\})
\end{verbatim}  
  \end{exampleblock}
\end{xframe}

%%%%%%%

\end{document}
